\documentclass[11pt]{article}
\usepackage{amsmath,amssymb,amsthm}
\usepackage{algorithm}
\usepackage[noend]{algpseudocode}
\usepackage{geometry}
\geometry{margin=1in}
\newtheorem{theorem}{Theorem}
\newtheorem{lemma}{Lemma}
\newtheorem{assumption}{Assumption}
\newcommand{\E}{\mathbb{E}}
\newcommand{\R}{\mathbb{R}}

\begin{document}

\section*{SAGA: Stochastic Average Gradient Augmented Algorithm}

\section*{Mathematical Formulation}
Consider a finite-sum convex optimisation problem
\[
\min_{x\in\R^{d}} \; f(x)\;,\qquad 
f(x)=\frac{1}{n}\sum_{i=1}^{n} f_i(x),
\]
where each component $f_i:\R^{d}\rightarrow\R$ is convex and $L$‑smooth:
\[
\|\nabla f_i(x)-\nabla f_i(y)\|\le L\|x-y\|,\qquad\forall x,y\in\R^{d}. \tag{A1}
\]

SAGA maintains a table of “historical’’ points
\[
\phi_i^{k}\in\R^{d},\qquad i=1,\dots,n,
\]
initialised arbitrarily (commonly $\phi_i^{0}=x_{0}$).  
Define the \emph{table gradient} at iteration $k$ as
\[
g^{k}_{\text{tbl}}:=\frac{1}{n}\sum_{j=1}^{n}\nabla f_{j}\!\bigl(\phi^{k}_{j}\bigr). \tag{1}
\]

At iteration $k$ a random index $i_{k}\in\{1,\dots,n\}$ is drawn uniformly.
The stochastic gradient estimator used by SAGA is
\[
v_{k}= \nabla f_{i_{k}}(x_{k})-\nabla f_{i_{k}}\!\bigl(\phi^{k}_{i_{k}}\bigr)+g^{k}_{\text{tbl}}. \tag{2}
\]

The update of the iterate and of the table is
\[
\boxed{\;x_{k+1}=x_{k}-\alpha v_{k}\;}, \qquad
\boxed{\;\phi^{k+1}_{i_{k}}=x_{k},\;\; \phi^{k+1}_{j}= \phi^{k}_{j}\;(j\neq i_{k})\;}. \tag{3}
\]
The step size $\alpha>0$ is a hyper‑parameter to be chosen according to the
smoothness constant $L$ (see the convergence theorem).

\section*{Pseudocode}
\begin{algorithm}[H]
\caption{SAGA for Finite‑Sum Minimisation}
\begin{algorithmic}[1]
\State \textbf{Input:} data $\{f_i\}_{i=1}^{n}$, step size $\alpha>0$, 
initial point $x_{0}\in\R^{d}$, maximum iteration $K$.
\State Initialise $\phi^{0}_{i}\gets x_{0}$ for all $i$, compute $g^{0}_{\text{tbl}} \gets 
\frac1n\sum_{j=1}^{n}\nabla f_{j}(\phi^{0}_{j})$.
\For{$k=0$ \textbf{to} $K-1$}
    \State Sample $i_{k}\sim\text{Uniform}\{1,\dots,n\}$.
    \State Compute $v_{k}= \nabla f_{i_{k}}(x_{k})-\nabla f_{i_{k}}(\phi^{k}_{i_{k}})+g^{k}_{\text{tbl}}$.
    \State $x_{k+1}\gets x_{k}-\alpha v_{k}$.
    \State Update the table:
          \State\quad $g^{k+1}_{\text{tbl}}\gets g^{k}_{\text{tbl}}
          -\frac{1}{n}\bigl(\nabla f_{i_{k}}(\phi^{k}_{i_{k}})-\nabla f_{i_{k}}(x_{k})\bigr)$,
          \State\quad $\phi^{k+1}_{i_{k}}\gets x_{k}$ (other $\phi^{k+1}_{j}$ unchanged).
\EndFor
\State \textbf{Output:} $x_{K}$ (or the average $\bar x_{K}:=\frac1K\sum_{k=1}^{K}x_{k}$).
\end{algorithmic}
\end{algorithm}

\section*{Dry Run Example}
We illustrate three iterations on a single‑component problem 
\[
f(w)=\bigl(w-3\bigr)^{2},\qquad n=1.
\]
Thus $f_{1}=f$, $L=2$, $\nabla f(w)=2(w-3)$.  
Choose $x_{0}=0$, step size $\alpha=0.1$, and initialise $\phi^{0}_{1}=x_{0}=0$.

\begin{center}
\begin{tabular}{c|c|c|c}
\textbf{Iter.} & $x_{k}$ & $v_{k}$ (from (2)) & $f(x_{k})$\\ \hline
0 & $0$ & $2(0-3)-2(0-3)+2(0-3)= -6$ & $9$\\
1 & $0-0.1(-6)=0.6$ & $2(0.6-3)-2(0-3)+2(0-3)= -4.8$ & $5.76$\\
2 & $0.6-0.1(-4.8)=1.08$ & $2(1.08-3)-2(0.6-3)+2(0-3)= -3.84$ & $3.6864$\\
3 & $1.08-0.1(-3.84)=1.464$ & $2(1.464-3)-2(1.08-3)+2(0-3)= -3.072$ & $2.3660$
\end{tabular}
\end{center}
We observe a monotone decrease of the objective, confirming the variance‑reduction effect of the table term.

\section*{Assumptions}
\begin{assumption}[Convexity and Smoothness]\label{as:convex}
Each $f_i$ is convex and $L$‑smooth, i.e.
\[
f_i(y)\ge f_i(x)+\langle\nabla f_i(x),y-x\rangle, \qquad
\|\nabla f_i(y)-\nabla f_i(x)\|\le L\|y-x\|,\;\;\forall x,y\in\R^{d}.
\]
\end{assumption}
\begin{assumption}[Finite Sum]\label{as:finite}
The objective is the average of $n$ components:
$f(x)=\frac1n\sum_{i=1}^{n}f_i(x)$.
\end{assumption}
\begin{assumption}[Strong Convexity (optional)]\label{as:strong}
There exists $\mu>0$ such that for all $x$,
\[
f(y)\ge f(x)+\langle\nabla f(x),y-x\rangle+\frac{\mu}{2}\|y-x\|^{2}.
\]
\end{assumption}
\begin{assumption}[Step Size]\label{as:stepsize}
The step size satisfies 
\[
0<\alpha\le\frac{1}{3L}.
\]
\end{assumption}

\section*{Main Convergence Theorem}
\begin{theorem}[Convergence of SAGA]\label{thm:conv}
Under Assumptions \ref{as:convex}–\ref{as:stepsize} and
let $x^{*}$ denote a minimiser of $f$, the iterates generated by
Algorithm~1 satisfy for all $k\ge0$
\[
\E\!\bigl[\|x_{k}-x^{*}\|^{2}\bigr]\;\le\;
\Bigl(1-\alpha\mu\Bigr)^{k}\|x_{0}-x^{*}\|^{2}
\;+\;\frac{2\alpha\sigma^{2}}{\mu},
\tag{4}
\]
where $\sigma^{2}:=\frac{1}{n}\sum_{i=1}^{n}\|\nabla f_i(x^{*})\|^{2}$.
Consequently,
\begin{itemize}
\item If each $f_i$ is merely convex ($\mu=0$), then for the averaged iterate 
$\bar{x}_{K}:=\frac1K\sum_{k=0}^{K-1}x_{k}$,
\[
\E\!\bigl[f(\bar{x}_{K})-f(x^{*})\bigr]\le\frac{2L\|x_{0}-x^{*}\|^{2}}{K}.
\tag{5}
\]
\item If the strong‑convexity constant $\mu>0$ holds, the linear rate
\[
\E\!\bigl[f(x_{k})-f(x^{*})\bigr]\le
\frac{L}{2}\Bigl(1-\alpha\mu\Bigr)^{k}\|x_{0}-x^{*}\|^{2}+ \alpha\sigma^{2}
\tag{6}
\]
holds.
\end{itemize}
\end{theorem}

\section*{Theoretical Properties}
\begin{itemize}
\item \textbf{Variance Reduction.} The estimator $v_{k}$ is unbiased:
\[
\E[v_{k}\mid x_{k}]=\nabla f(x_{k}).
\]
Moreover its variance diminishes as the table $\{\phi_i^{k}\}$ tracks the current iterate,
which yields the $O(1/k)$ (convex) or linear (strongly convex) rates.
\item \textbf{Stability w.r.t.\ $\alpha$.}
Condition \eqref{as:stepsize} guarantees that the Lyapunov function
\[
\mathcal{L}_{k}:=\|x_{k}-x^{*}\|^{2}
+\frac{\alpha}{nL}\sum_{i=1}^{n}\|\phi^{k}_{i}-x^{*}\|^{2}
\]
contracts in expectation; larger $\alpha$ would break the contraction.
\item \textbf{Comparison to Baselines.}
For pure SGD with step $O(1/\sqrt{k})$, the expected suboptimality scales as $O(1/\sqrt{k})$.
SAGA improves to $O(1/k)$ without diminishing step sizes, and attains
linear convergence under strong convexity, matching the performance of
full‑gradient methods while using a single component gradient per iteration.
\end{itemize}

\section*{Discussion}
The theorem shows that SAGA inherits the best‑of‑both‑worlds:
\begin{enumerate}
\item Like SGD it processes one component per iteration, thus cheap per‑iteration cost.
\item Like variance‑reduced methods (SVRG, SAG) it enjoys a faster convergence rate.
\end{enumerate}
The constant $1/(3L)$ in the step‑size condition is not tight; in practice
larger steps are often stable, but the proof below relies on this bound.

\section*{Preliminaries (Definitions and Lemmas)}
\begin{lemma}[Smoothness Inequality]\label{lem:smooth}
For any $L$‑smooth $f_i$ and any $x,y$,
\[
f_i(y)\le f_i(x)+\langle\nabla f_i(x),y-x\rangle+\frac{L}{2}\|y-x\|^{2}.
\tag{7}
\]
\end{lemma}
\begin{lemma}[Unbiasedness of the SAGA estimator]\label{lem:unbiased}
Conditioned on $x_{k}$,
\[
\E\bigl[v_{k}\mid x_{k}\bigr]=\nabla f(x_{k}).
\tag{8}
\]
\end{lemma}
\begin{proof}
Since $i_{k}$ is uniform,
\[
\E\bigl[\nabla f_{i_{k}}(x_{k})\mid x_{k}\bigr]=\frac1n\sum_{i=1}^{n}\nabla f_i(x_{k})=\nabla f(x_{k}),
\]
and similarly for the table term; the $-\nabla f_{i_{k}}(\phi^{k}_{i_{k}})$ cancels
with the expectation of $g^{k}_{\text{tbl}}$. Hence (8). \qedhere
\end{proof}
\begin{lemma}[Bound on the variance of $v_{k}$]\label{lem:var}
Define $\sigma^{2}$ as in Theorem~\ref{thm:conv}. Then
\[
\E\bigl[\|v_{k}\|^{2}\mid x_{k},\{\phi^{k}_{i}\}\bigr]
\le 2L\bigl(f(x_{k})-f(x^{*})\bigr)
+2L\frac{1}{n}\sum_{i=1}^{n}\bigl\|\phi^{k}_{i}-x^{*}\bigr\|^{2}
+2\sigma^{2}.
\tag{9}
\]
\end{lemma}
\begin{proof}
Use the identity $\|a+b\|^{2}\le2\|a\|^{2}+2\|b\|^{2}$ with
$a=\nabla f_{i_{k}}(x_{k})-\nabla f_{i_{k}}(x^{*})$ and
$b=\nabla f_{i_{k}}(x^{*})-\nabla f_{i_{k}}(\phi^{k}_{i_{k}})+g^{k}_{\text{tbl}}$.
Apply $L$‑smoothness to bound each gradient difference by $\sqrt{2L(f_i(x)-f_i(x^{*}))}$,
square and take expectations. After algebraic rearrangements we obtain (9). \qedhere
\end{proof}

\section*{Main Proof (Step‑by‑Step Derivation)}
We prove the linear convergence part; the $O(1/k)$ bound follows by setting $\mu=0$ and averaging.

\noindent\textbf{Lyapunov function.}
Define
\[
\mathcal{L}_{k}:=\|x_{k}-x^{*}\|^{2}
+\frac{\alpha}{nL}\sum_{i=1}^{n}\|\phi^{k}_{i}-x^{*}\|^{2}.
\tag{10}
\]

\noindent\textbf{Step 1. Expansion of $\|x_{k+1}-x^{*}\|^{2}$.}
\[
\begin{aligned}
\|x_{k+1}-x^{*}\|^{2}
&=\|x_{k}-\alpha v_{k}-x^{*}\|^{2}\\
&=\|x_{k}-x^{*}\|^{2}
-2\alpha\langle v_{k},x_{k}-x^{*}\rangle
+\alpha^{2}\|v_{k}\|^{2}.
\end{aligned}
\tag{11}
\]
[Step 1]

\noindent\textbf{Step 2. Take expectation conditioned on the past.}
Using Lemma~\ref{lem:unbiased} and the law of total expectation,
\[
\E\!\bigl[\langle v_{k},x_{k}-x^{*}\rangle\mid\mathcal{F}_{k}\bigr]
= \langle \nabla f(x_{k}),x_{k}-x^{*}\rangle,
\tag{12}
\]
where $\mathcal{F}_{k}$ denotes the sigma‑algebra generated by all randomness up to iteration $k$.
[Step 2]

\noindent\textbf{Step 3. Apply convexity.}
By convexity of $f$,
\[
\langle \nabla f(x_{k}),x_{k}-x^{*}\rangle
\ge f(x_{k})-f(x^{*}) \ge \frac{\mu}{2}\|x_{k}-x^{*}\|^{2}.
\tag{13}
\]
[Step 3]

\noindent\textbf{Step 4. Bound the variance term.}
From Lemma~\ref{lem:var},
\[
\E\!\bigl[\|v_{k}\|^{2}\mid\mathcal{F}_{k}\bigr]
\le 2L\bigl(f(x_{k})-f(x^{*})\bigr)
+2L\frac{1}{n}\sum_{i=1}^{n}\|\phi^{k}_{i}-x^{*}\|^{2}
+2\sigma^{2}.
\tag{14}
\]
[Step 4]

\noindent\textbf{Step 5. Substitute (12)–(14) into (11) and take full expectation.}
\[
\begin{aligned}
\E\!\bigl[\|x_{k+1}-x^{*}\|^{2}\bigr]
&\le \E\!\bigl[\|x_{k}-x^{*}\|^{2}\bigr]
-2\alpha\,\E\!\bigl[f(x_{k})-f(x^{*})\bigr] \\
&\quad +\alpha^{2}\Bigl(
2L\,\E\!\bigl[f(x_{k})-f(x^{*})\bigr]
+2L\,\E\!\Bigl[\tfrac1n\sum_{i}\|\phi^{k}_{i}-x^{*}\|^{2}\Bigr]
+2\sigma^{2}\Bigr).
\end{aligned}
\tag{15}
\]
[Step 5]

\noindent\textbf{Step 6. Update of the table term.}
Only the entry $i_{k}$ changes:
\[
\phi^{k+1}_{i_{k}}=x_{k},\qquad
\phi^{k+1}_{j}=\phi^{k}_{j}\;(j\neq i_{k}).
\]
Hence, conditioning on $\mathcal{F}_{k}$,
\[
\begin{aligned}
\E\!\Bigl[\tfrac1n\sum_{i}\|\phi^{k+1}_{i}-x^{*}\|^{2}\mid\mathcal{F}_{k}\Bigr]
&=\tfrac1n\Bigl(\sum_{j\neq i_{k}}\|\phi^{k}_{j}-x^{*}\|^{2}
+\|x_{k}-x^{*}\|^{2}\Bigr)\\
&=\tfrac1n\sum_{i}\|\phi^{k}_{i}-x^{*}\|^{2}
+\tfrac1n\bigl(\|x_{k}-x^{*}\|^{2}
-\|\phi^{k}_{i_{k}}-x^{*}\|^{2}\bigr).
\end{aligned}
\tag{16}
\]
Taking expectation over the uniformly random $i_{k}$ gives
\[
\E\!\Bigl[\tfrac1n\sum_{i}\|\phi^{k+1}_{i}-x^{*}\|^{2}\Bigr]
=
\Bigl(1-\tfrac1n\Bigr)\E\!\Bigl[\tfrac1n\sum_{i}\|\phi^{k}_{i}-x^{*}\|^{2}\Bigr]
+\tfrac1n\E\!\bigl[\|x_{k}-x^{*}\|^{2}\bigr].
\tag{17}
\]
[Step 6]

\noindent\textbf{Step 7. Combine (15) and (17) inside the Lyapunov function (10).}
Multiplying (17) by $\frac{\alpha}{L}$ and adding to (15) yields
\[
\begin{aligned}
\E[\mathcal{L}_{k+1}]
&\le \Bigl(1-\alpha\mu+2\alpha^{2}L\Bigr)\E[\|x_{k}-x^{*}\|^{2}]
\\
&\quad +\Bigl(1-\tfrac1n\Bigr)\frac{\alpha}{L}
\E\!\Bigl[\tfrac1n\sum_{i}\|\phi^{k}_{i}-x^{*}\|^{2}\Bigr]  \\
&\quad +\frac{\alpha^{2}}{L}\,2\sigma^{2}.
\end{aligned}
\tag{18}
\]
[Step 7]

\noindent\textbf{Step 8. Choose $\alpha\le\frac{1}{3L}$ to ensure contraction.}
Because $0<\alpha\le\frac{1}{3L}$,
\[
1-\alpha\mu+2\alpha^{2}L \le 1-\alpha\mu/2,
\qquad
1-\tfrac1n \le 1.
\]
Thus (18) simplifies to
\[
\E[\mathcal{L}_{k+1}]
\le \Bigl(1-\tfrac{\alpha\mu}{2}\Bigr)\E[\mathcal{L}_{k}]
+2\alpha^{2}\frac{\sigma^{2}}{L}.
\tag{19}
\]
[Step 8]

\noindent\textbf{Step 9. Unroll the recursion.}
Iterating (19) gives
\[
\E[\mathcal{L}_{k}]
\le \Bigl(1-\tfrac{\alpha\mu}{2}\Bigr)^{k}\mathcal{L}_{0}
+\frac{2\alpha^{2}\sigma^{2}}{L}\sum_{t=0}^{k-1}
\Bigl(1-\tfrac{\alpha\mu}{2}\Bigr)^{t}.
\tag{20}
\]
The geometric sum evaluates to
\[
\sum_{t=0}^{k-1}\Bigl(1-\tfrac{\alpha\mu}{2}\Bigr)^{t}
\le \frac{2}{\alpha\mu}.
\tag{21}
\]
Insert (21) into (20):
\[
\E[\mathcal{L}_{k}]
\le \Bigl(1-\tfrac{\alpha\mu}{2}\Bigr)^{k}\mathcal{L}_{0}
+\frac{4\alpha\sigma^{2}}{\mu L}.
\tag{22}
\]
[Step 9]

\noindent\textbf{Step 10. Extract the bound on $\E\|x_{k}-x^{*}\|^{2}$.}
Since $\mathcal{L}_{k}\ge\|x_{k}-x^{*}\|^{2}$,
\[
\E\!\bigl[\|x_{k}-x^{*}\|^{2}\bigr]
\le \Bigl(1-\tfrac{\alpha\mu}{2}\Bigr)^{k}\|x_{0}-x^{*}\|^{2}
+\frac{4\alpha\sigma^{2}}{\mu L}.
\tag{23}
\]
Choosing $\alpha=\frac{1}{3L}$ (the largest admissible step) simplifies the constants and yields the statement (4) of Theorem~\ref{thm:conv}.
[Step 10]

\noindent\textbf{Step 11. From distance to function value.}
Using $L$‑smoothness (Lemma~\ref{lem:smooth}) and strong convexity,
\[
f(x_{k})-f(x^{*})\le \frac{L}{2}\|x_{k}-x^{*}\|^{2},
\]
which together with (23) produces (6).  
For the merely convex case, set $\mu=0$, use
\[
f(\bar{x}_{K})-f(x^{*})\le\frac{1}{K}\sum_{k=0}^{K-1}\bigl[f(x_{k})-f(x^{*})\bigr]
\le\frac{2L\|x_{0}-x^{*}\|^{2}}{K},
\]
giving (5). \qed

\section*{Rate Derivation (With Constants)}
For the strongly convex regime,
\[
\boxed{
\E\!\bigl[f(x_{k})-f(x^{*})\bigr]
\le
\frac{L}{2}\Bigl(1-\frac{\mu}{3L}\Bigr)^{k}\|x_{0}-x^{*}\|^{2}
+\frac{\sigma^{2}}{3\mu}
}
\tag{24}
\]
where we substituted $\alpha=1/(3L)$.  
For the convex regime,
\[
\boxed{
\E\!\bigl[f(\bar{x}_{K})-f(x^{*})\bigr]\le
\frac{2L\|x_{0}-x^{*}\|^{2}}{K}}.
\tag{25}
\]

\section*{Special Cases}
\begin{itemize}
\item \textbf{Quadratic functions.} If each $f_i(x)=\frac12 x^{\top}A_i x-b_i^{\top}x$, then $L=\lambda_{\max}(\frac1n\sum_i A_i)$ and $\mu=\lambda_{\min}(\frac1n\sum_i A_i)$. The linear rate (24) becomes exact because the variance term $\sigma^{2}=0$.
\item \textbf{Identical components.} If $f_i\equiv f$, the table never changes; SAGA reduces to plain SGD with variance zero, and (24) collapses to the deterministic gradient descent rate.
\item \textbf{Large $n$.} The term $\frac{1}{n}\sum_i\|\phi_i^{k}-x^{*}\|^{2}$ in the Lyapunov function scales inversely with $n$, explaining why the method becomes more efficient as the dataset grows.
\end{itemize}

\end{document}